%-------------------------
% Career Description in Latex
% Original Author : Sourabh Bajaj
% Adaptation : Hyunggi Chang
% Draft content : Changyeob Lee
% Target : CERT / Incident Response / Security Monitoring
% File : career_description_no_metrics_ko.tex
% License : MIT
%------------------------

\documentclass[letterpaper,11pt]{article}

\usepackage{latexsym}
\usepackage[empty]{fullpage}
\usepackage{titlesec}
\usepackage{marvosym}
\usepackage[usenames,dvipsnames]{color}
\usepackage{verbatim}
\usepackage{enumitem}
\usepackage[hidelinks]{hyperref}
\usepackage{fancyhdr}
\usepackage[english]{babel}
\usepackage{tabularx}
\usepackage{amsmath}
\usepackage{kotex}

\pagestyle{fancy}
\fancyhf{}
\fancyfoot{}
\renewcommand{\headrulewidth}{0pt}
\renewcommand{\footrulewidth}{0pt}

\addtolength{\oddsidemargin}{-0.5in}
\addtolength{\evensidemargin}{-0.5in}
\addtolength{\textwidth}{1in}
\addtolength{\topmargin}{-0.5in}
\addtolength{\textheight}{1.0in}

\urlstyle{same}
\raggedbottom
\raggedright
\setlength{\tabcolsep}{0in}

\titleformat{\section}{
  \vspace{-4pt}\scshape\raggedright\large
}{}{0em}{}[\color{black}\titlerule \vspace{-2pt}]

\newcommand{\resumeItem}[1]{
  \item\small{{#1 \vspace{-2pt}}}
}

\newcommand{\resumeSubheading}[4]{
  \vspace{-1pt}\item
    \begin{tabular*}{0.97\textwidth}[t]{l@{\extracolsep{\fill}}r}
      \textbf{#1} & #2 \\
      \textit{\small#3} & \textit{\small #4} \\
    \end{tabular*}\vspace{-5pt}
}

\newcommand{\resumeProject}[2]{
  \vspace{-1pt}\item
    \begin{minipage}[t]{0.97\textwidth}
      \textbf{#1} \\
      \small{#2}
    \end{minipage}\vspace{-5pt}
}

\newcommand{\resumeTalk}[2]{
  \vspace{-1pt}\item
    \begin{minipage}[t]{0.97\textwidth}
      \textbf{#1} \\
      \textit{\small #2}
    \end{minipage}\vspace{-5pt}
}

\newcommand{\resumeSkill}[2]{
  \item\small{\textbf{#1}: #2}
}

\newcommand{\resumeAward}[3]{
  \item\small{\textbf{#1} / #2 / #3 \vspace{-2pt}}
}

\renewcommand{\labelitemii}{$\circ$}
\newcommand{\resumeSubHeadingListStart}{\begin{itemize}[leftmargin=*]}
\newcommand{\resumeSubHeadingListEnd}{\end{itemize}}
% 내용을 덜어내지 않고 페이지 수를 줄이기 위해 항목 사이 '단락 간격'만 없앱니다.
% resumeItem 이 59개라 항목당 몇 pt 만 줄여도 한 페이지가 확보됩니다.
% itemsep 은 건드리지 않습니다 — resumeItem 안에 이미 \vspace{-2pt} 가 있어 함께 줄이면 글자가 겹칩니다.
\newcommand{\resumeItemListStart}{\begin{itemize}[topsep=2pt,parsep=0pt,partopsep=0pt]}
\newcommand{\resumeItemListEnd}{\end{itemize}\vspace{-5pt}}

\begin{document}

%----------HEADING-----------------
\begin{tabular*}{\textwidth}{l@{\extracolsep{\fill}}r}
  \textbf{\href{https://info.yeopeva.me/}{\Large 이창엽 (Changyeob Lee)}} & Github: \href{https://github.com/YeoPEVA}{YeoPEVA} $|$ LinkedIn: \href{https://www.linkedin.com/in/changyeob-lee-yeopeva}{changyeob-lee-yeopeva} \\
  \href{https://info.yeopeva.me/}{https://info.yeopeva.me/} & Email: \href{mailto:lcy606@naver.com}{lcy606@naver.com} \\
  \href{https://noperfectsecurity.tistory.com/}{noperfectsecurity.tistory.com} & KST / GMT+9
\end{tabular*}

%-----------SUMMARY-----------------
\section{요약}
  \resumeSubHeadingListStart
    \resumeItem{CERT, 침해사고 대응, 보안 관제 직무를 목표로 보안 운영, 디지털 포렌식, 악성코드 분석, 보안 자동화 경험을 쌓아온 정보보호 전공 지원자입니다.}
    \resumeItem{공군 정보보호병으로 실운영 보안 환경에서 정보보호 장비 관리, 보안 관제, 보안 정책 및 차단 관리, 훈련 대응 업무를 수행하며 보안 운영 절차와 상황 공유 흐름을 경험했습니다.}
    \resumeItem{디지털 포렌식·침해대응 경진대회와 연구 활동을 통해 Windows Artifact, MFT, Registry, Prefetch, Event Log 기반 사고 타임라인 재구성, 증거 분석, 보고서 작성 경험을 쌓았습니다.}
    \resumeItem{악성코드와 익스플로잇 킷 분석 활동을 통해 파일 구조, 실행 흐름, 네트워크 행위, IoC 정리, YARA 룰 작성 관점의 분석 역량을 키웠습니다.}
    \resumeItem{MFT Parser 개선과 Stellar-Agent 프로젝트를 통해 분석 도구 개선, IaC 보안 정책 감사, SBOM 및 보고서 생성 자동화 경험을 보유하고 있습니다.}
    \resumeItem{부산소프트웨어마이스터고에서 2023년부터 전공동아리 멘토로 활동하고 있으며, 2026년에는 산학겸임교사로 정보보호관리 교과목을 담당하며 보안 지식을 교육 내용으로 구조화하고 전달하는 경험을 쌓았습니다.}
  \resumeSubHeadingListEnd

%-----------EXPERIENCE-----------------
\section{주요 경험}
  \resumeSubHeadingListStart

    \resumeSubheading
      {HSPACE FORENSIC LAB}
      {대한민국}
      {Member / HDFLAB}
      {2025.10 -- 현재}
      \resumeItemListStart
        \resumeItem{디지털 포렌식 및 침해사고대응 분야의 연구 커뮤니티 활동을 이어가고 있습니다.}
        \resumeItem{소비팟 소속으로 상·하반기마다 제공되는 문제를 풀이한 뒤 보고서로 작성하고 있습니다.}
        \resumeItem{분석 과정에서 확인한 내용을 팀 내 공유 문서로 정리하며 포렌식 분석 결과를 설명 가능한 형태로 구조화하는 경험을 쌓았습니다.}
      \resumeItemListEnd

    \resumeSubheading
      {SCH 사이버보안연구센터}
      {순천향대학교}
      {Member / 악성코드, 익스플로잇 킷 분석 및 연구}
      {2025.09 -- 현재}
      \resumeItemListStart
        \resumeItem{악성코드와 익스플로잇 킷 분석을 중심으로 보안 연구 활동을 수행하고 있습니다.}
        \resumeItem{랜섬웨어·인포스틸러를 주로 분석하며, 악성코드 분석·수집 자동화를 함께 연구하고 있습니다.}
        \resumeItem{파일 구조, 실행 흐름, 난독화 여부, 네트워크 행위, 지속성 확보 가능성을 중심으로 분석하며, 분석 결과를 IoC, 행위 요약, 탐지 가능 지점, 후속 분석 항목으로 정리하고 있습니다.}
      \resumeItemListEnd

    \resumeSubheading
      {부산소프트웨어마이스터고}
      {부산, 대한민국}
      {전공동아리 멘토 / 산학겸임교사 · 정보보호관리 교과목 담당}
      {2023.04 -- 현재}
      \resumeItemListStart
        \resumeItem{2023 -- 2026학년도 전공동아리 멘토로 연속 임명되어 학생 프로젝트를 지도하고 있습니다. 담당 팀은 2023·2024 마루, 2025 솔빗, 2026 다솜입니다.}
        \resumeItem{2026.03.09 -- 2026.07.16 산학겸임교사로 정보보호관리 교과목을 담당했습니다.}
        \resumeItem{정보보호관리 교과 내용을 학생들이 이해할 수 있는 수업 흐름으로 구조화하고, 보안 운영과 관리 관점의 개념을 교육 현장에 맞게 전달했습니다.}
        \resumeItem{이 활동을 통해 보안 지식을 설명 가능한 형태로 정리하는 역량과 교육 및 커뮤니케이션 역량을 함께 쌓았습니다.}
      \resumeItemListEnd

    \resumeSubheading
      {i-Keeper}
      {대구가톨릭대학교}
      {CERT팀 / DEV팀 Member}
      {2023.09 -- 2024.07}
      \resumeItemListStart
        \resumeItem{analyzeMFT 기반 MFT Parser 분석 및 개선 프로젝트를 수행했습니다.}
        \resumeItem{Windows 파일시스템 아티팩트 분석 관점에서 MFT 구조와 파싱 결과를 검토하고, 분석 결과 확인 편의성을 높이는 방향으로 도구 개선에 참여했습니다.}
        \resumeItem{DEV팀 장비부장으로 활동하며 장비 관리와 개발 지원 업무를 병행했습니다.}
      \resumeItemListEnd

    \resumeSubheading
      {Anti-root}
      {영남권}
      {팀장 및 설립자 / 정보보호 연구팀 운영}
      {2017.05 -- 2023.02}
      \resumeItemListStart
        \resumeItem{영남권 정보보호 연구팀을 설립·운영하며 자체 세미나, 프로젝트, 발표 활동을 주관했습니다.}
        \resumeItem{YARA, 이메일 포렌식, 허니팟, USB 복구, CloudGoat 등 다양한 보안 주제로 발표 및 실습을 진행했습니다.}
        \resumeItem{팀장으로서 주제 선정, 일정 관리, 발표 피드백, 결과물 정리를 담당했으며, 보안 학습 내용을 팀 내에서 공유 가능한 자료로 구조화했습니다.}
      \resumeItemListEnd

    \resumeSubheading
      {대한민국 공군}
      {대한민국}
      {정보보호병 825기 / 보안 장비 관리, 관제, 정책·차단 관리}
      {2021.04 -- 2023.01}
      \resumeItemListStart
        \resumeItem{실운영 보안 환경에서 정보보호 장비 관리, 보안 관제, 보안 정책·차단 관리 업무를 수행했습니다.}
        \resumeItem{보안 이벤트와 장비 상태를 확인하고, 정책·차단 요청 발생 시 적용 대상과 영향 가능성을 검토했습니다.}
        \resumeItem{운영 이슈 발생 시 담당자 간 상황 공유와 후속 보고 흐름을 경험했으며, 훈련 대응 과정에서 절차 기반 대응과 상황 전파의 중요성을 익혔습니다.}
        \resumeItem{반복적으로 발생하는 운영 이슈는 인수인계 자료와 운영 가이드 형태로 정리하여 업무 연속성을 높이는 데 기여했습니다.}
        \resumeItem{제44회 공군 정보통신 경연대회 해킹방어부문 최우수상(공군참모총장)을 수상했습니다.}
      \resumeItemListEnd

  \resumeSubHeadingListEnd

%-----------PROJECTS-----------------
\section{핵심 프로젝트 및 활동}
  \resumeSubHeadingListStart

    \resumeProject
      {악성코드 분석 및 YARA 룰 작성 경험}
      {2025.09 -- 현재 / 악성코드 샘플 분석, IoC 정리, YARA 룰 작성, 분석 보고서 작성}
      \resumeItemListStart
        \resumeItem{악성코드 샘플에 대해 정적·동적 분석을 수행하며 파일 구조와 행위 기반 흔적을 확인했습니다.}
      \resumeItemListEnd

    \resumeProject
      {Stellar-Agent / IaC-Rule Set \& MCP 기반 정책 감사 도구}
      {2025.03.03 -- 2025.12.18 / IaC Rule Set, MCP 기반 AI, 보안 정책 감사, SBOM 및 보고서 생성 자동화}
      \resumeItemListStart
        \resumeItem{IaC 분석 파트와 2025 한국데이터사이언스학회 동계 종합학술대회 논문 작성에 참여했습니다.}
        \resumeItem{IaC Rule Set을 기준으로 보안 정책 위반 가능성을 점검하는 감사 흐름 구현에 참여했습니다.}
        \resumeItem{MCP 기반 AI를 활용해 감사 결과를 해석하고 보고서 형태로 정리하는 기능 구현에 기여했으며, SBOM 및 분석 보고서 생성 흐름을 구성하여 반복 점검 결과를 문서화 가능한 형태로 구조화했습니다.}
        \resumeItem{2025 SCHU AI \& SW Festival SW 프로젝트 경진대회 대상(순천향대학교 SW중심대학사업단장)을 수상했습니다.}
      \resumeItemListEnd

    \resumeProject
      {금융보안아카데미 2025 사이버 위협 대응·분석 과정}
      {2025.07.08 -- 2025.11.07 / 금융권 위협 분석·대응 절차 학습 및 DFIR 중심 CTF 수행}
      \resumeItemListStart
        \resumeItem{금융권 사이버 위협 대응 절차와 분석 보고 흐름을 학습했습니다.}
        \resumeItem{DFIR 중심 CTF 문제를 분석하며 사고 원인 추적, 증거 분석, 타임라인 정리, 결과 문서화 과정을 수행했습니다.}
        \resumeItem{분석 결과를 발표자료와 보고서 형태로 정리하며 기술적 근거와 결론을 분리해 전달하는 경험을 쌓았습니다.}
        \resumeItem{금융보안아카데미 2025 최우수상(금융보안원장)을 수상했습니다.}
      \resumeItemListEnd

    \resumeProject
      {디지털 포렌식 · 침해대응 경진대회 분석 경험}
      {2022.12 -- 2025.12 / 디지털 증거 분석, 사고 타임라인 재구성, 분석 보고서 작성}
      \resumeItemListStart
        \resumeItem{디지털 포렌식·침해대응 경진대회에 지속적으로 참여하며 로그, 파일 아티팩트, 시스템 이벤트를 기반으로 침해사고 흐름을 분석했습니다.}
        \resumeItem{MFT, Registry, Prefetch, Event Log 등 복수 아티팩트를 교차 검증하여 공격자의 유입 경로, 실행 흔적, 내부 행위, 피해 범위를 추정했습니다.}
        \resumeItem{분석 결과는 사고 개요, 주요 증거, 타임라인, 원인 추정, 영향 범위, 후속 확인 항목 중심의 보고서로 정리했습니다.}
        \resumeItem{2025 제6회 호남 사이버보안 콘퍼런스 침해대응 경진대회 우수상, 2024 제5회 호남사이버보안컨퍼런스 대학간 침해대응/분석 경진대회 장려상, 제8--11회 디지털 범인을 찾아라 경진대회 장려상(4회 연속)을 수상했습니다.}
      \resumeItemListEnd

    \resumeProject
      {i-Keeper / MFT Parser 분석 및 개선 프로젝트}
      {2023.09 -- 2024.02 / Python, analyzeMFT 기반 MFT 도구 분석, GUI 및 기능 개선}
      \resumeItemListStart
        \resumeItem{analyzeMFT 기반 MFT Parser의 구조와 파싱 결과를 검토했습니다.}
        \resumeItem{MFT 분석 결과 확인 편의성을 높이고 파일시스템 기반 타임라인 재구성 과정을 보조하는 방향으로 도구 개선에 참여했습니다.}
      \resumeItemListEnd

    \resumeProject
      {BoB - 중고거래사기방지플랫폼 구현}
      {2020.09 -- 2020.12 / Python, PHP, 중고거래 사기 게시물 탐지 및 사용자 알림 플랫폼 구현}
      \resumeItemListStart
        \resumeItem{게시글 크롤링과 백엔드 구현을 담당했습니다.}
        \resumeItem{중고거래 게시글 데이터를 기반으로 사기 탐지와 사용자 알림 흐름을 구현했습니다.}
        \resumeItem{프로젝트 결과를 논문으로 확장하여 KCI 등재 논문과 한국정보보호학회 동계학술대회 발표로 연결했습니다.}
        \resumeItem{BoB 9기 디지털포렌식 트랙에서 TOP 14 자문단 평가까지 진출했습니다.}
      \resumeItemListEnd

  \resumeSubHeadingListEnd

%-----------SKILLS-----------------
\section{기술 및 관심 분야}
  \resumeSubHeadingListStart
    \resumeSkill{디지털 포렌식 · 침해대응 (DFIR)}{Windows·Linux 주요 아티팩트 분석, 타임라인 재구성, 침해사고 분석, 분석 보고서 작성}
    \resumeSkill{악성코드 분석 · 탐지 룰}{정적·동적 분석, PE 구조 분석, 파일·레지스트리·프로세스·네트워크 행위 분석, IoC 정리, YARA}
    \resumeSkill{보안 운영}{보안 관제, 보안 장비 관리, 정책·차단 관리, 훈련 대응, 상황 전파, 인수인계 문서화}
    \resumeSkill{엔지니어링}{Python, Backend, DevSecOps, IaC Rule Set, MCP 기반 AI, SBOM, 보고서 생성 자동화, AWS, OpenStack}
    \resumeSkill{관심 분야}{CTI, DevSecOps, 백엔드}
  \resumeSubHeadingListEnd

%-----------AWARDS-----------------
\section{수상 및 성과}
  \resumeSubHeadingListStart
    \resumeAward{장려상(한국포렌식학회장)}{제8--11회 디지털 범인을 찾아라 경진대회 (4회 연속)}{2022.12.15 / 2023.12.18 / 2024.12.17 / 2025.12.16}
    \resumeAward{최우수상(금융보안원장)}{금융보안아카데미 2025}{2025.11.07}
    \resumeAward{대상(순천향대학교 SW중심대학사업단장)}{2025 SCHU AI \& SW Festival SW 프로젝트 경진대회}{2025.11.05}
    \resumeAward{우수상(한국주택금융공사 사장)}{제1회 영남권 사이버 공격 방어 대회 대학팀 부문}{2025.10.15}
    \resumeAward{우수상(한국인터넷진흥원장)}{2025 제6회 호남 사이버보안 콘퍼런스 침해대응 경진대회}{2025.09.18}
    \resumeAward{최우수팀(대구가톨릭대학교 대학일자리플러스센터장)}{2024-2 창의도전학기 사례발표회}{2024.12.05}
    \resumeAward{우수상(대구가톨릭대학교 SW중심대학사업단장)}{2024 제22회 TOPCIT 정기평가}{2024.11.11}
    \resumeAward{장려상(한국정보보호학회 호남지부장)}{2024 제5회 호남사이버보안컨퍼런스 대학간 침해대응/분석 경진대회}{2024.10.02}
    \resumeAward{대상(경북ICT융합산업진흥협회)}{한-베트남 제2회 국제SW코딩 경진대회}{2024.06.18}
    \resumeAward{최우수상(공군참모총장)}{제44회 공군 정보통신 경연대회 해킹방어부문}{2022.10.26}
    \resumeAward{장려상(한국디지털포렌식학회장)}{2020년도 한국디지털포렌식학회 디지털포렌식 챌린지}{2020.11.20}
  \resumeSubHeadingListEnd

%-----------TALKS AND PAPERS-----------------
\section{발표 및 논문}
  \resumeSubHeadingListStart
    \resumeTalk
      {Stellar-Agent: IaC 룰셋 및 모델 컨텍스트 프로토콜(MCP)을 통합한 AI 에이전트 기반 하이브리드 DevSecOps 정책 감사 프레임워크}
      {2025 한국데이터사이언스학회 동계 종합학술대회 논문집, 2025.12.18, 논문번호 0032, p.118}
      \resumeItemListStart
        \resumeItem{정준영, 이창엽, 이정민, 이유식}
        \resumeItem{IaC 분석 파트와 논문 작성에 참여했습니다.}
      \resumeItemListEnd

    \resumeTalk
      {온라인 직거래 환경에서의 사기 피해에 대한 사전 예방 필요성 제언: 사기 판별 요소를 중심으로}
      {KCI 등재 논문, 경찰학연구 21(1), pp.249--272, 2021.03, DOI: 10.22816/polsci.2021.21.1.009}
      \resumeItemListStart
        \resumeItem{현장호, 임다연, 이창엽, 이진우, 이현호, 최원영}
        \resumeItem{BoB 프로젝트 결과를 논문으로 확장하여 온라인 직거래 사기 판별 요소와 예방 필요성을 정리했습니다.}
      \resumeItemListEnd

    \resumeTalk
      {중고 거래 사기 방지 플랫폼 제안}
      {한국정보보호학회 동계학술대회, 2020}
      \resumeItemListStart
        \resumeItem{현장호, 이현호, 임다연, 이진우, 이창엽, 오동빈, 최원영}
      \resumeItemListEnd

    \resumeTalk
      {국내 APT 및 기업 보안 사고에 대한 고찰}
      {한국정보보호학회 동계학술대회, 2019}
      \resumeItemListStart
        \resumeItem{이창엽, 이진우, 이정훈, 노무승, 박현상}
        \resumeItem{국내 APT 및 기업 보안 사고 사례를 조사하고 보안 사고 대응 관점에서 시사점을 정리했습니다.}
      \resumeItemListEnd

    \resumeTalk
      {카드해킹의 문제점과 대비방안}
      {한국정보보호학회 동계학술대회, 2018}
      \resumeItemListStart
        \resumeItem{이진우, 이정훈, 이창엽, 노무승, 오현수}
      \resumeItemListEnd

    \resumeTalk
      {RFID Skimming 해킹 기법에 관한 고찰}
      {한국정보보호학회 동계학술대회, 2017}
      \resumeItemListStart
        \resumeItem{김건우, 김태현, 성유원, 여상준, 이진우, 이창엽, 이달원}
      \resumeItemListEnd
  \resumeSubHeadingListEnd

%-----------CERTIFICATIONS-----------------
\section{자격 및 취득 이력}
  \resumeSubHeadingListStart
    \resumeItem{AWS Certified Solutions Architect - Associate (2023.07.26 -- 2026.07.26, 만료)}
    \resumeItem{디지털포렌식 2급 (2022.12.22 -- 2028.12.22)}
    \resumeItem{정보기기운용기능사 (2022.10 취득)}
    \resumeItem{Exterro ACE / AccessData (2022.01.11 -- 2024.01.11, 만료)}
    \resumeItem{인터넷정보관리사 2급 (2020.10 취득)}
    \resumeItem{리눅스마스터 2급 (2019.10 취득)}
    \resumeItem{네트워크관리사 2급 (2016.12 -- 2026.12.12)}
  \resumeSubHeadingListEnd

\end{document}
